\begin{center}
\textsc{Gauss} an \textsc{Gerling}.\quad Göttingen, 5.~September 1823.
\end{center}

\par\noindent\hspace*{\parindent}\mbox{.\ .\ .\ .\ .\ .}\quad Unbedingt wünsche ich, dass Sie den Heliotrop [auf dem
Inselsberg] genau in das Alignement mit dem Theodolithenplatz stellen und das Zelt
so viel wie möglich symmetrisch aufschlagen. Könnten Sie der dem Brocken
zugekehrten Seite eine dunkle Farbe geben, so wäre es so viel besser. Ich
\textls[120]{hoffe}, dass der dortige Heliotrop das Zelt weit überschreien soll; allein ich
habe 1822 doch den Fall gehabt, dass das \textls[120]{weisse} Zelt in Scharnhorst meine
Beobachtungen in Garssen einen ganzen Tag unbrauchbar machte, obwohl
dabei zu bemerken, dass erstens das Zelt mir seine Südseite, von der Sonne
beleuchtet, zukehrte, zweitens das Heliotroplicht sehr geschwächt war, welches
beides hier wegfällt. Inzwischen möchte ich doch aus Vorsicht den Heliotrop
\textls[120]{nicht excentrisch} aufgestellt wissen, da sonst das Zelt wohl ein wenig
beitragen könnte, den Winkel zu verfälschen. Ich vermuthe, wenn z.~B.\ das
Zelt \(\frac{1}{40}\) so viel Glanz hat wie der Heliotrop, und dieser so weit absteht, dass
die Distanz \(4^{\prime\prime}\) gross sein solle, der Winkel um \(0\mathord{\ooalign{\hfil\(\vphantom{0}^{\prime\prime}\)\hfil\cr\hfil\({,}\)\hfil\cr}}1\) verfälscht werden wird.\\*
\mbox{.\ .\ .\ .\ .\ .}\par
